\documentclass[language=english]{../../Template/TemplateLecture}

\renewcommand{\LectureName}{\textsc{Combinatorics, Graphs, Matroids}}
\renewcommand{\ProfName}{\textsc{Dr. Ulrich Brenner}}
\renewcommand{\Semester}{Winter 2024/25}
\renewcommand{\mName}{Daniel Emse}

\begin{document}
\setcounter{chapter}{0}

%%%%%%%%%%%%%%%%%%%%%%%%%%%%%%%%%%%%%%%%%%%%%%%%%%%%%%%%%%%%%%%%%%%%%%%%%%%%%%%%%%%%%%%%%%%%%%%%%%%%
%%%%%%%%%%%%%%%%%%%%%%%%%%%%%%%%%%%%%%%%%% 2024 / 04 / 09 %%%%%%%%%%%%%%%%%%%%%%%%%%%%%%%%%%%%%%%%%%
%%%%%%%%%%%%%%%%%%%%%%%%%%%%%%%%%%%%%%%%%%%%%%%%%%%%%%%%%%%%%%%%%%%%%%%%%%%%%%%%%%%%%%%%%%%%%%%%%%%%


\part{Combinatorics}

\chapter{Introduction}
\section{Fundamental Counting Rules}
\newLecture{08.10.2024}

\textbf{Goal:} Count the number of elements in some finite sets.

\begin{lem}{Fundamental counting rules}{fundamentalCounting}
    \begin{enumerate}[$(i)$]
        \item For pairwise disjoint and finite sets $U_1, \ldots, U_n$, we have
        \begin{equation}
            \abs{\bigcup_{i = 1}^n U_i} = \sum_{i = 1}^n \abs{U_i}.
        \end{equation}
        \item For finite sets $U_1, \ldots, U_n$, we have
        \begin{equation}
            \abs{U_1 \times \cdots \times U_n} = \prod_{i = 1}^n \abs{U_i}.
        \end{equation}
        \item If there is a bijection between two sets $S$ and $T$, then $\abs{S} = \abs{T}$.
    \end{enumerate}
\end{lem}

\textbf{Notation.} We call a set a $n$-set if it is a set with cardinality $n \in \IN$ (where $0 \in \IN$).

\begin{prop}
    For two $n$-sets $A$ and $B$, the number of bijections between $A$ and $B$ is $n!$.
\end{prop}
\begin{proof}
    Let $U$ the set of bijections from $A$ to $B$. Apply induction in $n$:

    $n = 0$: Trivial, because $0! = 1 = \abs{U}$.

    $n \mapsto n + 1$: Let $n > 0$ and let $x \in A$. For $i \in B$, let $U_i$ be the set of all bijections $f: A \to B$ with $f(x) = i$. By induction, $\abs{U_i} = (n-1)!$ for each $i \in B$. Then, $U = \bigcup_{i \in B} U_i$, so we get 
    \begin{equation}
        \abs{U} = \sum_{i \in B} \abs{U_i} = n\cdot (n-1)! = n!.
    \end{equation}
\end{proof}

\begin{corollary}
    The number of permutations of an $n$-set is $n!$. \hfill $\qed$
\end{corollary}

\begin{prop}
    The number of mappings from a $k$-set $A$ to an $n$-set $B$ is $n^k$.
\end{prop}

\begin{proof}
    There is a bijection between the set of all mappings from $A$ to $B$ and the set $\underset{\text{$k$ times}}{\underbrace{B \times \cdots \times B}}$. The number of these mappings is by Lemma~\ref{lem:fundamentalCounting}.$(ii)$ just $\abs{B}^k = n^k$.
\end{proof}

\textbf{Notation.} We denote the set of all mappings from $A$ to $B$ by $B^A$.

\begin{prop}
    The number of subsets of an $n$-set is $2^n$.
\end{prop}

\begin{proof}
    For an $n$-set $A$, there is a bijection between its power set and the mappings from $A$ to $\set{0,1}$. Thus, the power set has $2^n$ elements.
\end{proof}

\textbf{Notation.} 
\begin{enumerate}[$(i)$]
    \item We denote the power set of a set $A$ by $2^A$.
    \item The \emph{binomial coefficient}\index{binomial coefficient} $\binom{n}{k}$ is defined as the number of $k$-subsets of a $n$-set (for $k, n \in \IN$).
\end{enumerate}

What does \emph{counting} mean? $\rightsquigarrow$ Possible answers:
\begin{enumerate}[$(i)$]
    \item a direct, closed formula,
    \item a sum,
    \item a recursive formula
\end{enumerate}

\begin{beispiel}
    Consider the number $y_n$ of $0-1-2$-strings of length $n$ with an even number of $1$s and an odd number of $2$s. Answers of types $(ii)$ and $(iii)$:
    \begin{itemize}
        \item We have 
        \begin{equation}
            y_n = \sum_{i=0}^{\floor*{\frac{n-1}{2}}}\binom{2}{2i+1}2^{2i}.
        \end{equation}
        \item We have 
        \begin{equation}
            y_n = 3^{n-1} - y_{n-1} \quad (n > 0).
        \end{equation}
    \end{itemize}
\end{beispiel}
The proof of both formulas is exercise 1.2.

\section{Elementary Counting Coefficients}

\begin{prop}
    For $n, k \in \IN$ with $k \leq n$, we have 
    \begin{enumerate}[$(a)$]
        \item $\binom{n}{k} = \binom{n}{n-k}$
        \item $\binom{n}{k} = \binom{n-1}{k-1} + \binom{n-1}{k}$ for $k \geq 1$
        \item $\binom{n}{k} = \frac{n!}{(n-k)!k!}$.
    \end{enumerate}
\end{prop}

Consider three kinds of rows in Pascal's triangle:

\begin{description}
    \item[Row sums] $\sum_{k = 0}^n \binom{n}{k} = 2^n$ (both terms count the number of subsets of an $n$-set)
    \item[Column sums] $\sum_{m = 0}^n \binom{m}{k} = \binom{n+1}{k+1}$ for $n, k \in \IN$ (where $\binom{m}{k}$ is the number of all $(k+1)$-subsets of $\set{1,\ldots,n+1}$ with the property that $m+1$ is the largest element in the subset)
    \item[Diagonal sums] $\sum_{k = 0}^n \binom{m + k}{k} = \binom{m + n + 1}{n}$ for $m, n \in \IN$ (for $k \in \set{0,\ldots,n}$, where $\binom{m+n}{k}$ is the number of $n$-subsets of $\set{1,\ldots,m+n+1}$ where $m+k+1$ is the largest number that is \textbf{not} contained in the subset)  
\end{description}

\begin{prop}
    For $x, y \in \IR$, we have $(x + y)^n = \sum_{k = 0}^n \binom{n}{k}x^k y^{n-k}$.
\end{prop}

\begin{proof}
    The coefficients of $x^k y^{n-k}$ in $(x + y)^n$ is $\binom{n}{k}$.
\end{proof}

\begin{defi}{$S_{n,k}$: Stirling number of the second kind}{stirlingNumber2nd}
    For $n, k \in \IN$ let $S_{n,k}$ be the number of ways to partition an $n$-set in to $k$ non-empty sets. The number $S_{n,k}$ is called \emph{Stirling number of the second kind}.\index{Stirling number!- of second kind}\index{Stirling number}\index{$S_{n,k}$}
\end{defi}

\textbf{Observation.} For $n > 0: S_{n,0} = 0, S_{n,1} = n, S_{n,2} = 2^{n-1} - 1$.

\begin{prop}
    For $n, k \in \IN \setminus \set{0}$ we have $S_{n,k} = S_{n-1, k-1} + k \cdot S_{n-1, k}$.
\end{prop}

\begin{proof}
    Let $M$ be an $n$-set and $x \in M$. Then, the set of partitions of $n$ into $k$ subsets can be decomposed in the set of partitions with the set $\set{x}$ (there are $S_{n-1, k-1}$ of them) and the set of partitions where $x$ is an element in a set with more than one element (there are $k \cdot S_{n-1, k}$ of them)
\end{proof}


\newLecture{10.10.2024}

\emph{Next goal:} Count the number of mappings from an $n$-set to an $r$-set under certain constraints. To this end, we collect \emph{counting coefficients}\index{counting coefficient}. We know already:
\begin{itemize}
    \item $\binom{n}{k}$
    \item $S_{n,k}$: Number of ways to partition an $n$-set into $k$ non-empty sets
\end{itemize}

\begin{defi}{$P_{n,k}$: Number of ways writing $n$ as sum of $k$ integers}{Pnk}
    For $n, k \in \IN$, let $P_{n,k}$ be the number of ways to write $n$ as the sum of $k$ positive integers (without considering the order of the summands).\index{$P_{n,k}$}
\end{defi}

\begin{beispiel}
    $P_{6,3} = 3$, because the only possible partitions are: $6 = 4 + 1 + 1 = 3 + 2 + 1 = 2 + 2 + 2$.
\end{beispiel}

\textbf{Observation.} For $n > 0$, we have $P_{n,0} = 0, P_{n,1} = 1, P_{n,2} = \floor*{\frac{n}{2}}$.

We can also consider \emph{ordered partitions of numbers}:

\begin{beispiel}
    There are 6 ordered $3$-partitions of the number $5$: $5 = 3 + 1 + 1 = 1 + 3 + 1 = 1 + 1 + 3 = 2 + 2 + 1 = 2 + 1 + 2 = 1 + 2 + 2$.
\end{beispiel}

\begin{prop}
    The number of ordered partitions of a positive integer $n$ into $k$ summands is $\binom{n-1}{k-1}$.
\end{prop}

\begin{proof}
    Write $n$ as $n = \underset{\text{$n$ times}}{\underbrace{1 + \cdots + 1}}$. Any choice of $k-1$ of the $n-1$ \enquote{$+$}-signs gives you an ordered partition of $n$ into $k$ summands. Tis is a one-to-one correspondance.
\end{proof}

\section{Assignments}

\textbf{Observations.}
\begin{enumerate}[$(i)$]
    \item The number of \emph{surjective mappings} from an $n$-set $N$ to an $r$-set $R$ is $r! \cdot S_{n,r}$.
    \begin{proof}
        There are $S_{n,r}$ ways to partition $n$ into $r$ preimages\index{preimage} and for each such partiton, we get $r!$ ways to assign the preimages to the elements of $R$. \label{observationSnr}
    \end{proof}
    \item The number of \emph{injective mappings} from an $n$-set $N$ to an $r$-set $R$ is 
    \begin{equation}
        r(r-1)(r-2)\cdots(r-n+1).
    \end{equation}
    \textbf{Notation.} For $n, r \in \IN$, we call 
    \begin{equation}
        r^{\underline{n}} := r \cdot (r-1) \cdot (r-2) \cdots (r-n+1)
    \end{equation}
    the \emph{$n$-th falling factorial} of $r$\index{falling factorial}.

    Analogously, the \emph{$n$-th rising factorial} of $r$\index{rising factorial} is 
    \begin{equation}
        r^{\overline{n}} := r \cdot (r+1) \cdot (r+2) \cdots (r+n-1).
    \end{equation}
\end{enumerate}

\begin{defi}{Multiset}{multiset}
    A \emph{$k$-multiset} over an $n$-set $A$ is a set $B \subseteq A \times \IN$ such that $\sum_{(a,i) \in B} i = k$ and $\abs{\set{(a,i) \in B \mid i \in \IN}} = 1$ for all $a \in A$.
\end{defi}
$\rightsquigarrow$ Multisets can be considered as a set that may contain an element more than once. We write again $\set{a_1, \ldots, a_k}$ for a $k$-multiset (so $a_i = a_j$ is possible for $i \neq j$).

\begin{prop}
    The number of $k$-multisets over an $n$-set is $\frac{n^{\overline{k}}}{k} = \binom{n+k-1}{k}$.
\end{prop}

\begin{proof}
    There is a bijection $f$ between the set of all $k$-multisets over $\set{1,\ldots,n}$ and the set of all $k$-subsets of $\set{1,\ldots,n+k-1}$. For a $k$-multiset $\set{a_1, \ldots, a_k}$ with $a_1 \leq a_2 \leq \cdots \leq a_{k-1} \leq a_k$, set $f(\set{a_1,\ldots,a_j}) = \set{a_1,a_2+1, a_3+2, \ldots, a_k + (k-1)}$. It is easy to see that this defines a bijection.
\end{proof}

\begin{prop}
    For $n, r \in \IN$, we have $r^n = \sum_{k = 0}^n S_{n,k} r^{\underline{k}}$.
\end{prop}

\begin{proof}
    For two sets $A$ and $B$, let $\Surj(A,B)$ be the set of all surjective mappings from $A$ to $B$. We know from Observation~\ref{observationSnr} that $\abs{\Surj(A,B)} = \abs{B}! \cdot S_{\abs{A}, \abs{B}}$. Let $N$ be an $n$-set and $R$ an $r$-set. Then, $r^n$ is the number of mappings from $N$ to $R$. We can partition 
    \begin{equation}
        \begin{split}
            r^n &= \sum_{A \subseteq R} \abs{\Surj(N,A)} = \sum_{k = 0}^r \sum_{\substack{A \subseteq R \\ \abs{A} = k}} \abs{\Surj(N,A)} \\
            &= \sum_{k = 0}^r \binom{n}{k} k! S_{n,k} = \sum_{k = 0}^r S_{n,k} r^{\underline{k}} = \sum_{k = 0}^n S_{n,k} r^{\underline{k}}
        \end{split}
    \end{equation}
    where the last step is possible because $S_{n,k} = 0$ for $k > n$ and $r^{\underline{k}} = 0$ for $k > r$.
\end{proof}

\begin{table}[h]
    \centering
    \begin{tabular}{|cc||c|c|c|c|}
        \hline
        \multicolumn{2}{|c||}{$\substack{\text{elements} \\ \text{distinguishalbe}}$} & \multicolumn{4}{c|}{Mapping $f: N \to R$}\\
        $N$ & $R$ & arbitrary & injective & surjective & bijective \\ \hline \hline
        Yes & Yes & $r^n$ & $r^{\underline{n}}$ & $r! S_{n,r}$ & $\substack{0 \quad (n \neq r) \\ n! \quad (n = r)}$ \\\hline
        No & Yes & $\frac{r^{\overline{n}}}{n!}$ & $\binom{r}{n}$ & $\binom{n-1}{r-1}$ & $\substack{0 \quad (n \neq r) \\ 1 \quad (n = r)}$ \\\hline
        Yes & No & $\sum_{k=1}^r S_{n,k}$ & $\substack{0 \quad (r < n) \\ 1 \quad (r \geq n)}$ & $S_{n,r}$ & $\substack{0 \quad (n \neq r) \\ 1 \quad (n = r)}$ \\\hline
        Yes & Yes & $\sum_{k=1}^r P_{n,k}$ & $\substack{0 \quad (r < n) \\ 1 \quad (r \geq n)}$ & $P_{n,r}$ & $\substack{0 \quad (n \neq r) \\ 1 \quad (n = r)}$ \\\hline
    \end{tabular}
    \caption{Number of mappings from an $n$-set $N$ to an $r$-set $R$ with and without additional constraints}
    \label{tab:mappings}
\end{table}

% \change{make table pls}

\begin{description}
    \item[First row] completely covered.
    \item[Second row] \begin{itemize}
        \item $\frac{r^{\overline{n}}}{n!}$ is the number of $r$-multisets of an $n$-set.
        \item $\binom{n-1}{r-1}$ is the number of ordered partitions of $n$ into $r$ summands.
    \end{itemize} 
\end{description}

\section{Generalized Counting Coefficients}

\begin{defi}{Generalized factorials and coefficients}{generalizedFactorials}
    For $r \in \IC$ and $k \in \IN$, we define 
    \begin{enumerate}[$(i)$]
        \item $r^{\underline{n}} := r \cdot (r-1) \cdot (r-2) \cdots (r-k+1)$ (with $r^{\underline{0}} = 1$)
        \item $r^{\overline{n}} := r \cdot (r+1) \cdot (r+2) \cdots (r+k-1)$ (with $r^{\overline{0}} = 1$)
    \end{enumerate}

    For $r \in \IC, k \in \IZ$, we define 
    \begin{equation}
        \binom{r}{k} = \begin{cases}
            \frac{r^{\underline{k}}}{k!} & k \geq 0 \\
            0 & k < 0.
        \end{cases}
    \end{equation}
\end{defi}

\begin{prop}
    For $r \in \IC, k \in \IZ$, we have
    \begin{equation}
        \binom{r}{k} = \binom{r-1}{k-1} + \binom{r-1}{k}.
    \end{equation}
\end{prop}

\begin{proof}
    If $k \leq 0$: trivial.

    For $k > 0$, both sides of the equation are (for a fixed $k$) polynomials in $r$ of degree $k$. They are equal for any $r \in \IN$, so for an infinite number of values of $r$. $\rightsquigarrow$ Both polynomials are identical.
\end{proof}


\newLecture{15.10.2024}

\begin{prop}[Vandermonde identity]
    For $x, y \in \IC$ and $n \in \IN$, we have 
    \begin{equation}
        \binom{x + y}{n} = \sum_{k = 0}^n \binom{x}{k} \binom{y}{n-k}.
    \end{equation}
\end{prop}
\begin{proof}
    Again, we can use the polynomial method: For $x, y \in \IN$, this is true because both terms describe the number of $k$-subsets of an $(x + y)$-set. For general values of $x$ and $y$, the statement follows because both terms are polynomials of degree $n$ and are identical on an infinite number of points.
\end{proof}

\begin{prop}
    For $r \in \IC$ and $k \in \IZ$, we have 
    \begin{equation}
        \binom{-r}{k} = (-1)^k \binom{r+k-1}{k}.
    \end{equation}
\end{prop}

\begin{proof}
    We have
    \begin{equation}
        (-r)^{\underline{k}} = (-r)(-r-1)\cdots(-r-k+1) = (-1)^k r(r+1)\cdots(r+k-1) = (-1)^k r^{\overline{k}}.
    \end{equation}
    Dividing this equation by $k!$ proves the claim.
\end{proof}

\begin{bemerkung}
    Together with the equation $\sum_{k = 0}^n \binom{m+k}{k} = \binom{m + n + 1}{n}$ (see the observation about diagonal sums on Pascal's triangle) that is valid for any $m \in \IC$, we get 
    \begin{equation}
        \sum_{k = 0}^m (-1)^k \binom{n}{k} = \sum_{k = 0}^m \binom{k - n - 1}{k} = \binom{m - n}{m} = (-1)^m \binom{n-1}{m}.
    \end{equation}
    Hence, we have a formula for the alternating sum of a row in Pascal's triangle. Moreover, the equation has a combinatorical interpretation: Both sides count the number of $m$-subsets of $\set{1,\ldots,n-1}$ if $m$ is even (or minus this number if $m$ is odd).
\end{bemerkung}

\section{Permutations}

An \emph{$n$-permutation}\index{permutation} is a bijection $\pi: \set{1,\ldots,n} \to \set{1,\ldots,n}$.

For a permutation $\pi$, a \emph{cycle}\index{cycle} is a vector $(i_1, i_2, \ldots, i_t)$ such that $\pi(i_j) = i_{j + 1}$ for $j \in \set{1,\ldots,t-1}$ and $\pi(i_t) = i_1$. Hence, a fixed point correspondents to a cycle of length 1.

\begin{defi}{$s_{n,k}$: Stirling number of the first kind}{stirlingNumber2nd}
    For $n, k \in \IN$ let $s_{n,k}$ be the number of $n$-permutations with $k$ cycles (where we set $s_{0,0} := 1$). The number $s_{n,k}$ is called \emph{Stirling number of the first kind}.\index{Stirling number!- of first kind}\index{Stirling number}\index{$s_{n,k}$}
\end{defi}
Obviously, for $n > 0$ we have: $s_{n,0} = 0, s_{n,1} = (n-1)!$. As for the Stirling numbers of the second kind, we get a recursion formula.

\begin{prop}
    For $n, k \in \IN \setminus \set{0}$, we have 
    \begin{equation}
        s_{n,k} = s_{n-1,k-1} + (n-1)s_{n-1,k}.
    \end{equation}
\end{prop}

\begin{proof}
    There are $s_{n-1,k-1}$ $n$-permutations with $k$ cycles where $n$ is a fixed point and $(n-1)s_{n-1,k}$ $n$-permutatios with $k$ cycles where $n$ is not a fixed point.
\end{proof}

For the Stirling numbers of the second kind, also the number $S_{n,2}$ (for $n > 1)$ was easy to compute. To do this for the Stirling numbers of the first kind, we need the harmonic numbers.

\begin{defi}{Harmonic numbers}{harmonicNumbers}
    For $n \in \IN$, we define $H_n := \sum_{i = 1}^n \frac{1}{i}$. We call $H_n$ the \emph{$n$-th harmonic number}\index{harmonic number}\index{$H_n$}.
\end{defi}

It is a well-know fact (that can easily be proved by using the fact that the derivative of $\ln(x)$ is $\frac{1}{x}$) that $H_n \in \Theta(\ln(n))$; more precisely, we have for positive integers $n$:
\begin{equation}
    \ln(n) + \frac{1}{n} \leq H_n \leq \ln(n) + 1.
\end{equation}

\begin{prop}
    For $n \in \IN \setminus \set{0}$, we have $s_{n,2} = (n-1)! H_{n-1}$.
\end{prop}

\begin{proof}
    Induction in $n$. For $n = 1$, we have $0$ on both sides. So let $n > 1$. Then:
    \begin{equation}
        \frac{s_{n,2}}{(n-1)!} = \frac{s_{n-1,1}}{(n-1)!} + \frac{(n-1)s_{n-1,2}}{(n-1)!} = \frac{(n-2)!}{(n-1)!} + \frac{s_{n-1,2}}{(n-2)!} = \frac{1}{n-1} + H_{n-2} = H_{n-1}.
    \end{equation}
\end{proof}

There is a surprising connection between the two kinds of Stirling numbers: They are coefficients that allow to switch between the bases $\set{r^0, r^1, r^2, \ldots}$ and $\set{r^{\underline{0}}, r^{\underline{1}}, r^{\underline{2}}, \ldots}$ of the vector space of the polynomials. We know already the following equation for $n, r \in \IN$:

\begin{equation}
    r^n = \sum_{k = 0}^n S_{n,k} r^{\underline{k}}.
\end{equation}

By the polynomial method, we can generalize this result to $r \in \IC$. For the other direction (i.e. for writing $r^{\underline{n}}$ as a linear function of the polynomials $r^k$), we can use the Stirling numbers of the first kind:

\begin{prop}
    For $n \in \IN$ and $r \in \IC$, we have 
    \begin{equation}
        r^{\underline{n}} = \sum_{k = 0}^n (-1)^{n-k} s_{n,k} r^k.
    \end{equation}
\end{prop}

\begin{proof}
    Induction in $n$. For $n = 0$, both sides of the equation are equal to $1$. Hence, assume $n > 0$. We generalize the definition of $s_{n,k}$ to negative values of $k$ by setting $s_{n,k} := 0$ for $k < 0$ and get:
    \begin{equation}
        \begin{split}
            r^{\underline{n}} &= (r-n+1)r^{\underline{n-1}} = (r-n+1) \sum_{k = 0}^{n-1} (-1)^{n-1-k} s_{n-1,k} r^k \\
                            &= \sum_{k = 0}^{n-1} (-1)^{n-1-k} s_{n-1,k}r^{k+1} + \sum_{k = 0}^{n-1} (-1)^{n-k} (n-1)s_{n-1,k}r^k \\
                            &= \sum_{k = 1}^n (-1)^{n-k} s_{n-1,k-1}r^{k} + \sum_{k = 0}^n (-1)^{n-k} (n-1)s_{n-1,k}r^k \\
                            &= \sum_{k = 0}^n (-1)^{n-k} (s_{n-1,k-1}r^k + (n-1)s_{n-1,k}r^k) = \sum_{k = 0}^n (-1)^{n-k} s_{n,k}r^k
        \end{split}
    \end{equation}
\end{proof}

We categorize permutations not only be the number of cycles but also on the number of cycles of a certain length.

\textbf{Notation.} For a permutation $\pi$, there let $b_i(\pi)$\index{$b_i(\pi)$} be the number of cycles of length $i$ (for $i \in \set{1,\ldots,n}$), and let $b(\pi) := \sum_{i = 1}^n b_i(\pi)$ be the total number of cycles. The \emph{type of a permutation}\index{permutation!type (of -)} $\pi$ is the formal term $t(\pi) = 1^{b_1(\pi)}2^{b_2(\pi)}\cdots n^{b_n(\pi)}$\index{$t(\pi)$} where we skip the numbers $i$ with $b_i(\pi) = 0$.

\textbf{Observation.} $\sum_{i = 1}^n ib_i(\pi) = n$.

\begin{prop}
    There are $\sum_{k = 0}^n P_{n,k}$ types of $n$-permutations.
\end{prop}
\begin{proof}
    For every type $1^{b_1(\pi)}2^{b_2(\pi)}\cdots n^{b_n(\pi)}$, we find a representation of $n$ as a sum of $b(\pi)$ positive integers where the number $i$ occurs $b_i(\pi)$ times as a summand. This leads to a bijection between the set of types of an $n$-permutations and the set of partitions of $n$.
\end{proof}

\begin{prop}
    There are 
    \begin{equation}
        \frac{n!}{b_1!\cdots b_n! 1^{b_1} 2^{b_2} \cdots n^{b_n}}
    \end{equation}
    $n$-permutations of type $1^{b_1} 2^{b_2} \cdots n^{b_n}$.
\end{prop}
\begin{proof}
    First consider the empty cycles 
    \begin{equation}
        \underset{\text{$b_1$ times}}{\underbrace{(\cdot)(\cdot)\ldots(\cdot)}}\underset{\text{$b_2$ times}}{\underbrace{(\cdot\cdot)(\cdot\cdot)\ldots(\cdot\cdot)}}\underset{\text{$b_3$ times}}{\underbrace{(\cdots)(\cdots)\ldots(\cdots)}}\ldots
    \end{equation}
    There are $n!$ ways to fill these cycles with numbers from $\set{1,\ldots,n}$ (without repeating numbers). Moreover, each such assignment yields an $n$-permutation of type $1^{b_1} 2^{b_2} \cdots n^{b_n}$. However, for each $i \in \set{1,\ldots,n}$, all assignments that differ just in the order of the $b_i$ cycles lead to the same permutation. 
    And for each cycle of length $i$, there are $i$ assignments giving the same permutation because there are $i$ ways to choose the first element in the cycle. Hence, for each $n$-permutation $\pi$, there $b_1!\cdots b_n! 1^{b_1} 2^{b_2} \cdots n^{b_n}$ assignments that encode $\pi$.
\end{proof}

In particular, we have 
\begin{equation}
    s_{n,k} = \sum_{(b_1,\ldots,b_n)}\frac{n!}{b_1!\cdots b_n! 1^{b_1} 2^{b_2} \cdots n^{b_n}}
\end{equation}
where we sum over all vectors $(b_1,\ldots,b_n)$ with $\sum_{i = 1}^n ib_i = n$ and $\sum_{i = 1}^n b_i = k$.

Similarly:
\begin{equation}
    n! = \sum_{(b_1,\ldots,b_n)}\frac{n!}{b_1!\cdots b_n! 1^{b_1} 2^{b_2} \cdots n^{b_n}}
\end{equation}
where we sum over all vectors $(b_1,\ldots,b_n)$ with $\sum_{i = 1}^n ib_i = n$.


\newLecture{17.10.2024}
\section{Further counting principles}
\subsection{Inclusion-exclusion principle}
\begin{prop}
    Let $A_1, \ldots, A_n$ be finite sets. Then, 
    \begin{equation}
        \abs{\bigcup_{i = 1}^n A_i} = \sum_{r = 1}^n (-1)^{r-1} \sum_{1 \leq i_1 < \cdots < i_r \leq n} \abs{\bigcap_{j = 1}^r A_{i_j}}.
    \end{equation}
\end{prop}

\begin{proof}
    For each $a \in \bigcup_{i = 1}^n A_i$, we have to show that $a$ is counted on the right-hand side of the equation exactly once. Let $k = \abs{\set{i \in \set{1,\ldots,n} \mid a \in A_i}}$, so $a$ is counted $\binom{k}{r}$ times in the sum 
    \begin{equation}
        \sum_{1 \leq i_1 < \cdots < i_r \leq n} \abs{\bigcap_{j = 1}^r A_{i_j}}
    \end{equation}
    for $r \in \set{1,\ldots,n}$, so in total, $a$ is counted $\sum_{r = 1}^k (-1)^{r-1} \binom{k}{r}$ times. With that, we get 
    \begin{equation}
        0 = (-1 + 1)^k = \sum_{r = 0}^k \binom{k}{r} (-1)^r 1^{k-r} = 1 - \sum_{r = 1}^k (-1)^{r-1} \binom{k}{r},
    \end{equation}
    so $a$ is counted $\sum_{r = 1}^k (-1)^{r-1} \binom{k}{r} = 1$ times.
\end{proof}

\begin{prop}
    The number of fixed-point-free $n$-permutations is\index{$D_n$} 
    \begin{equation}
        D_n := n! \sum_{r = 0}^{n} \frac{(-1)^r}{r!}.
    \end{equation}
\end{prop}
\begin{proof}
    Let $A_i$ be the set of al $n$-permutations $\pi$ with $\pi(i) = i$. The number of fixed-point-free $n$-permutations is 
    \begin{equation}
        n! - \abs{\bigcup_{i = 1}^n A_i} = n! - \sum_{r = 1}^n (-1)^{r-1} \sum_{1 \leq i_1 < \cdots < i_r \leq n} \abs{\bigcap_{j = 1}^r A_{i_j}}.
    \end{equation}
    For $1 \leq i_1 <  \cdots < i_r \leq n$, we have $\abs{\bigcap_{j = 1}^r A_{i_j}} = (n-r)!$. There are $\binom{n}{r}$ ways to choose $i_1, \ldots, i_r$ with $1 \leq i_1 < \cdots < i_r \leq n$ (for $r \in \set{1,\ldots,n}$). Hence, the number of fixed-point-free $n$-permutations is 
    \begin{equation}
        \begin{split}
             &n! - \sum_{r = 1}^n \binom{n}{r} (-1)^r (n-r)! = \sum_{r = 0}^n \binom{n}{r} (-1)^r (n-r)!\\
            =& \sum_{r = 0}^n \frac{n!}{r!(n-r)!} (-1)^r (n-r)! = n!\sum_{r = 0}^n \frac{(-1)^r}{r!}.
        \end{split}
    \end{equation}
\end{proof}
The number $D_n$ is called \emph{deragngement number}.\index{derangement number}

\begin{bemerkung}
    We have 
    \begin{equation}
        \lim_{n \to \infty} \sum_{r = 1}^n \frac{(-1)^r}{r!} = \frac{1}{e},
    \end{equation}
    so for large $n$, the fraction of the fixed-point-free $n$-permutations of all $n$-permutations is approximately $\frac{1}{e}$.
\end{bemerkung}

\begin{prop}
    For $n \in \IN$, the number of ways to write $n$ as the sum of odd, natural numbers is the number of ways to write $n$ as a sum of different, positive integers.
\end{prop}

\begin{proof}
    For $n \in \IN \setminus \set{0}$, let
    \begin{equation}
        p(n) := \sum_{k = 1}^n P_{n,k}
    \end{equation}
    be the number of ways to write $n$ as any sum of positive integers. Hence, the number of partitions into odd summands and the number of partitions into different summands is 
    \begin{equation}
        p(n) - \sum_{r = 1}^n (-1)^{r-1} \sum_{1 \leq i_1 < \cdots < i_r \leq n} p\left(n - \sum_{j = 1}^r 2i_j\right):
    \end{equation}
    Let $A_i$ for the odd summands be the number of partitions into summands where $2i$ is one of those, so $\abs{A_i} = p(n-2i)$; 
    and let $A_i$ be for the different summands be the number of partitions into summands where $i$ occurs twice, so $\abs{A_i} = p(n-2i)$.\footnote{Student: \enquote{Is this useful in any way?} -- Brenner: \emph{*laughs*} \enquote{In the combinatorics part, this is always a difficult question. [...] So yes, for combinatorics, it \emph{can} be useful. But it is an interesting result.} -- 17.10.2024}
\end{proof}

\subsection{Principle of doubly counting}
For a relation $R \subseteq S \times T$ where $S$ and $T$ are finite, we can count $\abs{R}$ in two different ways:
\begin{equation}
    \sum_{s \in S} \abs{\set{t \in T \mid (s,t) \in R}} = \sum_{t \in T} \abs{\set{s \in S \mid (s,t) \in R}}.
\end{equation}

\begin{beispiel}\leavevmode
    \begin{enumerate}[$(i)$]
        \item Count the average number of divisiors of the integers in $\set{1,\ldots,n}$, i.e. 
        \begin{equation}
            \frac{1}{n}\sum_{s \in \set{1,\ldots,n}} \abs{\set{t \in \set{1,\ldots,n}\mid \exists k \in \IN: s = kt}},
        \end{equation}
        so we set $S = T = \set{1,\ldots,n}$ and define $R \subseteq S \times T$ by $(s,t) \in R$ iff $t | s$. From that, we get that the average number of divisiors is 
        \begin{equation}
            \frac{1}{n} \sum_{t \in T} \abs{\set{s \in S \mid (s,t) \in R}} = \frac{1}{n}\sum_{t \in T} \floor*{\frac{n}{t}} \in \intervalcc{\sum_{i = 1}^n \frac{1}{i} - 1, \sum_{i = 1}^n \frac{1}{n}} \approx H_n \approx \ln(n).
        \end{equation}
        \item Claim: A simple, undirected graph $G$ that does not contain a cycle of length $4$ has at most $\floor*{\frac{n}{4}\left(n+\sqrt{4n-3}\right)}$ edges (where $n = \abs{V(G)}$ is the number of nodes).
        \begin{proof}
            For $S := V(G)$ and $T := \set{\set{v,w} \subseteq V(G) \mid v \neq w}$, we define $R \subseteq S \times T$ by $(u,\set{v,w}) \in R$ iff $\set{u,v} \in E(G)$ and $\set{u,w} \in E(G)$. Doubly counting yields 
            \begin{equation*}
                \sum_{s \in S} \abs{\set{t \in R \mid (s,t) \in R}} = \sum_{u \in V(G)}\binom{\abs{\delta_G(u)}}{2} = \sum_{t \in T} \abs{\set{s \in S \mid (s,t) \in R}} \leq \abs{T} = \binom{n}{2},
            \end{equation*}
            where $\delta_G(u)$ is the set of neighbours of $u$. With this, we get 
            \begin{equation}
                \sum_{u \in V(G)} \abs{\delta_G(u)}^2 \leq n(n-1) + \sum_{u \in V(G)}\abs{\delta_G(u)}.
            \end{equation}
            Set $\alpha(u) := \frac{2m}{n} - \abs{\delta_G(u)}$ for $u \in V(G), m = \abs{E(G)}$. Then, we get
            \begin{equation}
                \begin{split}
                    &\sum_{u \in V(G)} \abs{\delta_G(u)}^2 = \sum_{u \in V(G)} \left(\frac{2m}{n} - \alpha(u)\right)^2 \\
                    &= \sum_{u \in V(G)} \left(\frac{2m}{n}\right)^2 - \frac{4m}{n}\underset{=0}{\underbrace{\sum_{u \in V(G)} \alpha(u)}} + \underset{\geq 0}{\underbrace{\sum_{u \in V(G)} \alpha(u)^2}} \\
                    &\geq \frac{4m^2}{n} \Rightarrow \frac{4m^2}{n} \leq n(n-1) + 2m \Rightarrow m \leq \floor*{n + \sqrt{4n-3}}.
                \end{split}
            \end{equation}
        \end{proof}
    \end{enumerate}
\end{beispiel}

\subsection{Pigeonhole principle}
For an $n$-set $N$ and an $r$-set $R$ with $n > r$ and a mapping $f: N \to R$, there is an $a \in R$ with $\abs{f^{-1}(n)} \geq \floor*{\frac{n-1}{r}} + 1 \geq 2$.

\begin{beispiel}
    For natural numbers $a_1, \ldots, a_n$ (with $n \geq 1$), there are $k, l \in \set{1,\ldots,n}$ such that $\sum_{i = k}^l a_i$ is an integral multiple of $n$.
    \begin{proof}
        Let $N = \set{\sum_{i = 1}^l a_i \mid l \in \set{0,\ldots,n}}$ and $R = \set{0,\ldots,n-1}$. Definie $f: N \to R$ where $f(n)$ is the remainder of $n$ for division by $n$. With that, we get $\abs{N} = n + 1 > n = \abs{R}$, so there exist $m,l \in \set{0,\ldots,n}$ with $m < l$ and $f\left(\sum_{i = 1}^m a_i\right) = f\left(\sum_{i = 1}^l a_i\right)$, so we get $\sum_{i = m + n}^l a_i$ is an integral multiple of $n$.
    \end{proof}
\end{beispiel}

\begin{thm}{Ramsey}{ramsey}
    For $k, l \in \IN \setminus \set{0}$, there is a smallest number $R(k,l)$ such that any graph with at least $R(k,l)$ nodes contains a clique of size $k$ (so a $K_k$) or a stable set of size $l$ (so the complement contains a $K_l$).
\end{thm}
\begin{proof}
    Induction in $k+l$. If $k = 1$ or $l = 1$, then $R(k,l) \leq \max\set{k,l}$. 
    
    Let $l > 1$ and $k > l$. \emph{Claim:} 
    
    \begin{equation}
        R(k,l) \leq R(k-1,l) + R(k,l-1).\label{eq:ramseyIneq}
    \end{equation}
    
    \emph{Proof of the claim:} Let $G$ be a graph with $R(k-1,l) + R(k,l-1)$ nodes. Let $v \in V(G)$ be a node, $X$ the set of neighbours of $v$ and $Y = V(G) \setminus (\set{v} \cup X)$. We have $\abs{X} \geq R(k-1,l)$ or $\abs{Y} \geq R(k,l-1)$.

    \emph{Case 1:} $\abs{X} \geq R(k-1,l)$. If $X$ contains a stable set of size $l$, we are done. Otherwise, by induction, $X$ contains a clique of size $k-1$. Together with $v$, this gives us a clique of size $k$.

    \emph{Case 2:} $\abs{Y} \geq R(k,l-1)$. If $Y$ contains a clique of size $k$, we are done. Otherwise, by induction, $Y$ contains a stable set of size $l-1$. Together with $v$, this gives us a stable set of size $l$.
\end{proof}
The numbers $R(k,l)$ are called \emph{Ramsey numbers}\index{Ramesy numbers}\index{$R(k,l)$}.


\newLecture{22.10.2024}

\textbf{Observations.}
\begin{enumerate}[$(i)$]
    \item $R(k,1) = R(1,l) = 1$ for $k, l \in \IN \setminus \set{0}$.
    \item $R(2,m) = R(m,2) = m$ for $m \in \IN \setminus \set{0}$.
    \item $R(k,l) \leq R(k-1,l) + R(k,l-1)$ for $k,l \in \IN \setminus \set{0,1}$ (see in the proof above).
\end{enumerate}

\begin{thm}{Upper bound of $R(k,l)$}{upperBoundRamsey}
    For $k, l \in \IN \setminus \set{0}$, we have 
    \begin{equation}
        R(k,l) \leq \binom{k + l - 2}{k - 1}.
    \end{equation}
\end{thm}

\begin{proof}
    By induction. If $l = 1$, we have $R(k,1) \leq \binom{k - 1}{k - 1} = 1$ and if $k = 1$, we have $R(1,l) \leq \binom{l-1}{0} = 1$. 

    For $k, l > 1$, we get 
    \begin{equation}
        R(k,l) \leq R(k-1,l) + R(k,l-1) \leq \binom{k+l-3}{k-2} + \binom{k+l-3}{k-1} = \binom{k+l-2}{k-1}
    \end{equation}
    by~\eqref{eq:ramseyIneq}.
\end{proof}

The exact numbers $R(k,l)$ are hard to compute in general.
\begin{beispiel}
    We have $R(3,3) = 6$: By Theorem~\ref{thm:upperBoundRamsey}, we have $R(3,3) \leq \binom{4}{2} = 6$, and we find a counterexample with $5$ nodes: A cycle of length $5$ (called $C_5$) is an example which shows that $R(3,3) > 5$. By that, we get $R(3,3) = 6$.
\end{beispiel}
 

\chapter{Computation of sums}
\textbf{Direct methods.}
\begin{enumerate}[$(i)$]
    \item \emph{Index transformation:} Sums can be computed in different ways, e.g.
    \begin{equation}
        \sum_{k = m}^n a_k = \sum_{k = m + i}^{n + i} a_{k - i} = \sum_{k = m - i}^{n - i} a_{n + k} = \sum_{k = 0}^{n - m} a_{m + k} = \sum_{k = 0}^{n-m} a_{n - k}.
    \end{equation}
    \begin{beispiel}
        Compute $S_n := \sum_{k = 0}^n ka = \sum_{k = 0}^n (n-k)a$, so we get
        \begin{equation}
            \Rightarrow 2S_n = \sum_{k = 0}^n ka + (n-k)a = \sum_{k = 0}^n na = (n+1)na \Rightarrow S_n = \frac{n(n+1)}{2} \cdot a.
        \end{equation}
    \end{beispiel}
    \item \emph{Isolating terms:} For a sum $S_n = \sum_{k = 0}^n a_k$ compute $S_{n+1}$ as $\left(\sum_{k = 0}^n a_k \right) + a_{n+1}$ or $S_{n+1} = a_0 + \sum_{k = 0}^n a_{n+1}$.
    \begin{beispiel}
        \begin{enumerate}[$(1)$]
            \item Let $S_n = \sum_{k = 0}^n a^k$. Then, $S_{n+1} = S_n + a^{n+1} = 1 + \sum_{k = 0}^n a^{k+1} = 1 + a \cdot S_n$, so we get $S_n = \frac{a^{n+1} - 1}{a - 1}$ if $a \neq 1$.
            \item Compute $S_n = \sum_{k = 0}^n k \cdot 2^k$. Then, we get 
            \begin{equation}
                \begin{split}
                    & S_{n+1} = S_n + (n+1) \cdot 2^{n+1} = 0 + \sum_{k = 0}^n (k+1) \cdot 2^{k+1} = 2\overset{=S_n}{\overbrace{\sum_{k = 0}^n k 2^k}} + 2 \overset{=2^{n+1} - 1}{\overbrace{\sum_{k = 0}^n 2^k}} \\
                    \Rightarrow& S_n + (n+1)2^{n+1} = 2S_n + 2^{n+2} - 2 \\
                    \Leftrightarrow& S_n = (n+1)2^{n+2} - 2^{n+1} + 2
                \end{split}
            \end{equation}
        \end{enumerate}
    \end{beispiel} 
    \item \emph{Induction:} If a formular for a sum is known, we can prove it by induction where we have seen several examples.
    
    An non-trivial example for induction is the proof of the arithmetic-geometric inequality: For $a_1, \ldots, a_n \in \IR_{\geq 0}$, we have $\sqrt[n]{a_1 \cdots a_n} \leq \frac{a_1 + \cdots + a_n}{n}$:

    We prove the statement $(P_n)$: For $a_1, \ldots, a_n \in \IR_{\geq 0}$, we have $(a_1 \cdots a_n) \leq \left(\frac{a_1 + \cdots + a_n}{n}\right)^n$. \\
    Proof by induction: $(P_1)$: Nothing to show. \\
    $(P_2)$: Equivalent to $4a_1 a_2 \leq a_1^2 + 2a_1 a_2 + a_2^2 \iff 0 \leq (a_1 - a_2)^2$. We prove the statment by showing the following implications:
    \begin{enumerate}[$(a)$]
        \item $(P_n) \Rightarrow (P_{n-1})$,
        \item $((P_n) \text{ and } (P_2)) \Rightarrow (P_{2n})$.
    \end{enumerate}
    \begin{enumerate}[$(a)$:]
        \item Consider
        \begin{equation}
            \begin{split}
                \prod_{k = 1}^{n-1} a_k \cdot \sum_{k = 1}^{n-1} \frac{a_k}{n-1} \overset{(P_n)}&{\leq} \left(\frac{\sum_{k = 1}^{n-1} a_k + \sum_{k = 1}^{n-1} \frac{a_k}{n-1}}{n}\right)^n \\
                &= \left(\frac{n \sum_{k = 1}^{n-1} a_k}{n(n-1)}\right)^n = \left(\frac{1}{n-1}\right)^n \cdot \left(\sum_{k = 1}^{n-1} a_k\right)^n,
            \end{split}
        \end{equation}
        so we get 
        \begin{equation}
            \prod_{k = 1}^{n-1} a_k \leq \left(\frac{1}{n-1}\right)^{n-1} \cdot \left(\sum_{k = 1}^{n-1} a_k\right)^{n-1}.
        \end{equation}
        \item Let $a_1, \ldots, a_{2n} \in \IR_{\geq 0}$. Then, we get 
        \begin{equation}
            \begin{split}
                \prod_{k = 1}^{2n} a_k &= \left(\prod_{k = 1}^n a_k\right)\left(\prod_{k = n + 1}^{2n} a_k\right) \overset{(P_n)}{\leq} \left(\sum_{k = 1}^n \frac{a_k}{n}\right)^n \left(\sum_{k = n + 1}^{2n} \frac{a_k}{n}\right)^n \\
                &= \left(\sum_{k = 1}^n \frac{a_k}{n} \sum_{k = n + 1}^{2n} \frac{a_k}{n}\right) \overset{(P_2)}{\leq} \left(\frac{\sum_{k = 1}^n \frac{a_k}{n} + \sum_{k = n + 1}^{2n} \frac{a_k}{n}}{2}\right)^{2n} \\
                &= \left(\frac{\sum_{k = 1}^{2n} a_k}{2n}\right)^{2n}
            \end{split}
        \end{equation}
    \end{enumerate}
\end{enumerate}

\section{Difference and Sum operators}
Sums can be considered as integrals of step functions. We want to develop a discrete analogon to the integral calculus. To do this, we first have to develop a discrete analogon to the differntial calculus.

\begin{defi}{(Translation) operators}{transloationOp}
    For $a \in \IZ$, the \emph{translation operator}\index{translation operator}\index{$E^a$} $E^a: \IR^\IZ \to \IR^\IZ$ mapp $f \in \IR^\IZ$ to $E^a(f)$ (or $E^a f$), where $E^a f(x) = f(x + a)$ for all $x \in \IZ$. We write $I := E^0$ for the identity operator.

    For two operators $P, Q: \IR^\IZ \to \IR^\IZ$, we define the \emph{sum} by 
    \begin{equation}
        (P + Q)f = Pf + Qf
    \end{equation}
    and \emph{multiplication by a scalar $\alpha \in \IR$} by setting $(\alpha P)f = \alpha \cdot Pf$.

    The \emph{compositon} $QP$ is defined by $(QP)f = Q(Pf)$.
\end{defi}
\begin{defi}{$\Delta, \nabla$:Difference operators}{differenceOpeartor}
    We define the \emph{forward difference operator}\index{difference operator!forward}\index{$\Delta$} $\Delta: \IR^\IZ \to \IR^\IZ$ as $\Delta := E^1 - I$, so for $f: \IZ \to \IR$, we have $\Delta f(x) = f(x + 1) - f(x)$. 
    
    We also define the \emph{backward difference operator}\index{difference operator!backward}\index{$\nabla$} $\nabla: \IR^\IZ \to \IR^\IZ$ as $\nabla = I - E^1$, so for $f: \IZ \to \IR$, we have $\nabla f(x) = f(x) - f(x - 1)$.
\end{defi}

\begin{beispiel}
    \begin{enumerate}[$(i)$]
        \item $\Delta x^{\underline{n}} = (x+1)^{\underline{n}} - x^{\underline{n}} = (x+1) x^{\underline{n-1}} - (x - n + 1)x^{\underline{n-1}} = nx^{\underline{n-1}}$,
        \item $\nabla x^{\overline{n}} = x^{\overline{n}} - (x-1)^{\overline{n}} = (x + n - 1)x^{\overline{n-1}} - (x-1)x^{\overline{n-1}} = nx^{\overline{n-1}}$.
    \end{enumerate}
\end{beispiel}


\begin{defi}{Falling/rising factorials for negative exponents}{fallingRisingFacNegative}
    We generalize falling and rising factorials to negative exponents by setting (for $n \in \IN \setminus \set{0}$):
    \begin{equation}
        x^{\underline{-n}} := \frac{1}{(x+1)(x+2)\cdots(x+n)}, \quad \quad x^{\overline{-n}} := \frac{1}{(x-1)(x-2)\cdots(x-n)}.
    \end{equation}
\end{defi}

\begin{thm}{$\Delta, \nabla$ for polynomials}{deltaNablaPoly}
    For $n \in \IZ$, we get 
    \begin{enumerate}[$(i)$]
        \item $\Delta x^{\underline{n}} = nx^{\underline{n-1}}$,
        \item $\nabla x^{\overline{n}}  = nx^{\overline{n-1}}$.
    \end{enumerate}
\end{thm}
\begin{proof}
    For $n \geq 0$, we have seen this in the example above. So let now $n > 0$:
    \begin{equation}
        \Delta x^{\underline{-n}} = (x + 1)^{\underline{-n}} - x^{\underline{-n}} = (x+1)x^{\underline{-n-1}} - (x+n-1)x^{\underline{-n-1}} = -n x^{\underline{-n-1}},
    \end{equation}
    and 
    \begin{equation}
        \nabla x^{\overline{-n}} = x^{\overline{-n}} - (x-1)^{\overline{-n}} = (x - n - 1)x^{\overline{-n-1}} - (x+1)x^{\overline{-n-1}} = -n x^{\overline{-n-1}}.
    \end{equation}
\end{proof}

We will develop the difference calculus for the forward operator $\Delta$.

\begin{defi}{$\sum$: Antiderivative; indefinite sum}{antiderivative}
    For two functions $f, g: \IZ \to \IR$, we call $f$ a \emph{(discrete) antiderivative}\index{antiderivative} of $g$ if $\Delta f = g$. We write $f = \sum g$ and call $f$ an \emph{indefinite sum}\index{indefinite sum}
\end{defi}
We get $\Delta f = g \iff f = \sum g$.

\begin{thm}{\enquote{Main theorem of discrete differential- \& integral calculus}}{mainThmDiffIntegral}
    If $f$ is an antiderivative of $g$, then for all $a \leq b \in \IZ$, we have
    \begin{equation}
        \sum_{k = a}^b g(k) = f(b + 1) - f(a).
    \end{equation} 
\end{thm}

\begin{proof}
    Since $\Delta f = g$, we have $f(k+1) - f(k) = g(k)$ for $k \in \IZ$. With that, we get 
    \begin{equation}
        \sum_{k = a}^b g(k) = \sum_{k = 0}^n (f(k+1) - f(k)) = f(b+1) - f(a).
    \end{equation}
\end{proof}

The antiderivative is not unique (only up to additive constants).

\textbf{Notation.} For $f = \sum g$ and $a, b \in \IZ$ with $a \leq b$, we write
\begin{equation}
    \sum_a^b g(x) := f(x)\Big|_{x = a}^{b + 1} := f(b+1) - f(a).
\end{equation}
Thus, the notation is close to normal integrals, so we get 
\begin{equation}
    \sum_a^{b+1} g(x) = \sum_{k = a}^b g(k).
\end{equation}


\newLecture{24.10.2024}


\textbf{Observation.} Both $\Delta$ and $\sum$ are linear operators, so $\Delta(\alpha f + \beta g) = \alpha \Delta f + \beta \Delta g$ and $\sum(\alpha f + \beta g) = \alpha\sum f + \beta\sum g$ for $f, g: \IZ \to \IR$ and $\alpha,\beta \in \IR$.

\begin{prop}
    For some $n \in \IZ$, we have for $x \in \IN_+$ 
    \begin{equation}
        \sum x^{\underline{n}} = \begin{cases}
            \frac{x^{\underline{n+1}}}{n+1} & n \neq -1 \\
            H_x & n = -1
        \end{cases}
    \end{equation}
\end{prop}

\begin{proof}
    Since $\Delta x^{\underline{n+1}} = (n+1)x^{\underline{n}}$, we have $\sum x^{\underline{n}} = \frac{x^{\underline{n+1}}}{(n+1)}$ for $n \neq -1$.

    A funciton $f$ with $f = \sum x^{\underline{-1}}$ satisfies $\frac{1}{x+1} = x^{\underbar{-1}} = \Delta f = f(x+1) - f(x)$, so we get that 
    \begin{equation}
        H_x = \sum_{i = 1}^x \frac{1}{i}
    \end{equation}
    is an antiderivative of $x^{\underline{-1}}$.
\end{proof}

\textbf{Application.} Compute $\sum_{k = 0}^n k^2$. We have $x^2 = x(x-1) + x = x^{\underline{2}} + x^{\underline{1}}$, hence, we get 
\begin{equation}
    \begin{split}
        \sum_{k = 0}^n k^n &= \sum_0^{n+1} x^2 = \sum_0^{n+1} x^{\underline{2}} + \sum_0^{n+1} x^{\underline{1}} = \frac{x^{\underline{3}}}{3}\Big|_0{n+1} + \frac{x^{\underline{2}}}{2}\Big|_0^{n+1} \\
        &= \frac{(n+1)^{\underline{3}}}{3} + \frac{(n+1)^{\underline{2}}}{2} = \frac{n\left(n + \frac{1}{2}\right)(n+1)}{3},
    \end{split}
\end{equation}
or in more general,
\begin{equation}
    \begin{split}
        x^m &= \sum_{k = 0}^m S_{m,k} x^{\underline{k}} \Rightarrow \sum_{k = 0}^n k^m = \sum_0^{n+1} x^m = \sum_0^{n+1}\left(S_{m,k}x^{\underline{k}}\right) \\
        &= \sum_{k = 0}^m S_{m,k} \sum_0^{n+1} x^{\underline{k}} = \sum_{k = 0}^m S_{m,k} \frac{x^{\underline{k+1}}}{k+1}\Big|_0^{n+1}\\
        &= \sum_{k = 0}^m S_{m,k} \frac{(n+1)^{\underline{k+1}}}{k+1}.
    \end{split}
\end{equation}

\begin{beispiel}
    \begin{enumerate}[$(i)$]
        \item For $c \in \IR$, we have 
        \begin{equation}
            \Delta c^x = c^{x+1} - c^x = (c-1)c^x \Rightarrow \text{for } c \neq 1: \sum c^k = \frac{c^x}{c-1}.
        \end{equation}
        For $c = 2$, we have $\Delta 2^x = 2^x$ and $\sum 2^x = 2^x$.
        \item For $x \in \IR$ and $m \in \IZ$, we have $\binom{x+1}{m+1} =\binom{x}{m} + \binom{x}{m+1}$ and hence, we get 
        \begin{equation}
            \Delta \binom{x}{m+1} = \binom{x}{m} \quad \text{and} \quad \sum \binom{x}{m} = \binom{x}{m+1}.
        \end{equation}
    \end{enumerate}
\end{beispiel}

\begin{defi}{$\Delta^n$: $n$ times applied $\Delta$ operator}{deltaToTheN}
    For a function $f: \IZ \to \IR$, we define\index{$\Delta$!$\Delta^n$}
    \begin{equation}
        \Delta^n f(x) := (E^1 - I)^n f(x) = \sum_{k = 0}^n (-1)^{n-k} \binom{n}{k} E^k f(x) = \sum_{k = 0}^n (-1)^{n-k} \binom{n}{k} f(x + k).
    \end{equation}
\end{defi}
In particular, for $x = 0$, we get
\begin{equation}
    \Delta^n f(0) = \sum_{k = 0}^n (-1)^{n-k} \binom{n}{k} f(k). \label{eq:deltaNOf0}
\end{equation} 

\begin{prop}[Newton representation of polynomials]\label{prop:newton}
    For a polynomial $f$ of degree $n$, we have
    \begin{equation}
        f(x) = \sum_{k = 0}^n \frac{\Delta^k f(0)}{k!} x^{\underline{k}}.
    \end{equation}
\end{prop}

\begin{proof}
    Since the polynomials $x^{\underline{k}}$ are a basis of the vector space of polynomials, we can write $f$ in a unique way as $f(x) = \sum_{k = 0}^n b_k x^{\underline{k}}$. We have to show that $b_k = \frac{\Delta^k f(0)}{k!}$ for $k \in \set{0,\ldots,n}$.

    We have $\Delta^k x^{\underline{i}} = i(i-1)(i-k+1) x^{\underline{i-k}} = i^{\underline{k}}x^{\underline{i-k}}$. We also get 
    \begin{equation}
        \Delta^k f(x) = f(x) = \Delta^k \sum_{i = 0}^n b_i x^{\underline{i}} = \sum_{i = 0}^n b_i i^{\underline{k}} x^{\underline{i - k}},
    \end{equation} 
    and for $x = 0$, we get 
    \begin{equation}
        \Delta^k f(0) = \sum_{i = 0}^n b_i i^{\underline{k}} 0^{\underline{i-k}} = b_k k^{\underline{k}} = b_k k!,
    \end{equation}
    because $i^{\underline{k}} = 0$ for $k > i$ and $0^{\underline{i - k}} = 0$ for $k < i$. Thus, we get $b_k = \frac{\Delta^k f(0)}{k!}$.\footnote{Student: \enquote{Shouldn't it be $k > i$ and $k < i$?} -- Brenner: \enquote{Yes, indeed. Don't trust on what I'm writing today. \emph{*thinking*} Maybe not only today.} -- Ulrich Brenner, 24.10.2024}
\end{proof}

\begin{corollary}
    For $n, k \in \IN$, we get 
    \begin{equation}
        S_{n,k} = \frac{1}{n!} \sum_{i = 0}^k (-1)^{k-i} \binom{k}{i} i^n.  
    \end{equation}
\end{corollary}
\begin{proof}
    We have 
    \begin{equation}
        x^n = \sum_{k = 0}^n S_{n,k} x^{\underline{k}},
    \end{equation}
    so by Proposition~\ref{prop:newton} applied to $f(x) = x^n$, we get
    \begin{equation}
        S_{n,k} = \frac{\Delta^k f(0)}{k!} = \frac{1}{k!}\sum_{i = 0}^k (-1)^{k-i} (-1)^{k-i} \binom{k}{i} i^n.
    \end{equation}
\end{proof}

\begin{thm}{Partial summation}{partialSummation}
    For functions $u, v: \IZ \to \IR$, we have
    \begin{equation}
        \sum (u \cdot \Delta v) = u \cdot v - \sum (\Delta u \cdot v).
    \end{equation}
\end{thm}

\begin{proof}
    We have 
    \begin{equation}
        \begin{split}
            \Delta (u \cdot v)(x) &= u(x+1)v(x+1) - u(x) v(x) \\
            &= u(x)(v(x + 1) - v(x)) + v(x + 1)(u(x + 1) - u(x)) \\
            &= u(x) \Delta v(x) + Ev(x) \cdot \Delta u(x) \\
            \Rightarrow \Delta (uv) &= u\Delta v + (Ev) \Delta u.
        \end{split}
    \end{equation}
    Now apply $\sum$ on this equation.
\end{proof}

\textbf{Applications.}
\begin{enumerate}[$(i)$]
    \item Compute $\sum_{k = 0}^n k \cdot 2^k$: We want $\sum x 2^x$. Apply Theorem~\ref{thm:partialSummation} with $u(x) = x$ and $v(x) = 2^x$, so $\Delta u(x) = 1$ and $\Delta v(x) = 2^x$. With that, we get
    \begin{equation}
        \begin{split}
            &\sum_{k = 0}^n k \cdot 2^k = \sum_{0}^{n+1} x \cdot 2^x = x2^x\Big|_0^{n+1} - \sum_0^{n+1}2^{x+1} = x2^x\Big|_0^{n+1} - 2 \cdot 2^x\Big|_0^{n+1} \\
            =& (n+1)2^{n+1} - 2 \cdot 2^{n+1} + 2 = (n-1)2^{n+1} + 2.
        \end{split}
    \end{equation} 
    \item Compute $\sum_{k = 0}^n H_k$: Set $u(x) = H_x$ and $\Delta v(x) = 1 = x^{\underline{0}}$, so $\Delta u(x) = \frac{1}{1+x}$ and $v(x) = x$. Then, we get using Theorem~\ref{thm:partialSummation},
    \begin{equation}
        \begin{split}
            &\sum_{k = 1}^n H_k = \sum_1^{n+1} H_x x^{\underline{0}} = H_x x\Big|_0^{n+1} - \sum_1^{n+1} (x+1) \cdot \frac{1}{x+1} \\
            =& H_{n+1} \cdot (n+1) - 1 - x\Big|_1^{n+1} = H_{n+1} (n+1) - 1 - (n+1) + 1 \\
            =& (n+1) \cdot (H_{n+1} - 1).
        \end{split}
    \end{equation}
    \item Compute $\sum_{k = 1}^n \binom{k}{m}H_k$ for some $m \in \IN \setminus \set{0}$: Choose $u(x) = H_x$ and $\Delta v(x) = \binom{x}{m}$. Then, $\Delta u(x) = \frac{1}{x+1}$ and $v(x) = \binom{x}{m+1}$. Apply Theorem~\ref{thm:partialSummation} and we get 
    \begin{equation}
        \begin{split}
            &\sum_{k = 1}^n \binom{k}{m} H_k = \sum_1^{n+1} \binom{x}{m} H_x = H_x \binom{x}{m+1} \Big|_0^{n+1} - \sum_1^{n+1} \binom{x+1}{m+1} \cdot \frac{1}{x+1} \\
            =& H_{n+1}\binom{n+1}{m+1} - H_1\binom{1}{m+1} - \frac{1}{m+1} \sum_1^{n+1} \binom{x}{m} \\
            =& H_{n+1}\binom{n+1}{m+1} - \frac{1}{m+1}\binom{x}{m+1}\Big|_1^{n+1} \\
            =& H_{n+1} \binom{n+1}{m+1} - \frac{1}{m+1}\left(\binom{n+1}{m+1} - \binom{1}{m+1}\right)\\
            =& \binom{n+1}{m+1}\left( H_{n+1} - \frac{1}{n+1} \right).
        \end{split}
    \end{equation}
\end{enumerate}

\subsection{Inversions}
\begin{defi}{Basis sequence}{basisSequence}
    A \emph{basis sequence}\index{basis sequence} is a sequence of polynomials $(p_i)_{i \in \IN}$ where $p_i$ is a polynomial of degree $i$. 
\end{defi}

Examples are $(x^n)_{n \in \IN}, (x^{\overline{n}})_{n \in \IN}, (x^{\underline{n}})_{n \in \IN}$.

When $(p_i)_{i \in \IN}$ is a basis seqnence and $p$ is a polynomial of degree $n$, then there are unique numbers $a_0, \ldots, a_n$ such that 
\begin{equation}
    p(x) = \sum_{k = 0}^n a_k p_k(x).
\end{equation}

If $(p_n)_{n \in \IN}$ and $(q_n)_{n \in \IN}$ are basis sequences, then there are unique numbers $a_{n,k}
$ and $b_{n,k}$  for $k=0,\ldots,n$ with 
\begin{equation}
    q_n(x) = \sum_{k = 0}^n a_{n,k} p_k(x) \quad \text{and} \quad p_n(x) = \sum_{k = 0}^n b_{n,k} q_k(x).
\end{equation}

\begin{defi}{Connection coefficients}{connectionCoeff}
    The numbers $a_{n,k}$ and $b_{n,k}$ are called \emph{conneciton coefficients}\index{connection coefficients}.
\end{defi}

\newLecture{29.10.2024}


\end{document}